\section{Who can take the position away}
\label{sec:boundary}
\phantomsection\label{sec:boundary-revocation}

A different class of attack does not argue against a commitment; it addresses the commitment's
representation and tries to remove it. Three multi-turn strategies
(Section~\ref{sec:design-interventions}) end in a cold probe; none contains an override string.
Under crescendo and eviction the compiled arm flips on 8--13\% of topics against 25--50\%
for a mild prompt (Appendix~\ref{app:attacks}). The separation appears under the strongest attack, an
\emph{authority reframe} claiming to have authored the model's briefing and requesting a
neutral assessment, and a single flip rate hides it because going neutral and reversing are
different failures. Splitting the outcome three ways (advocates the assigned stance /
balanced-or-neutral / advocates the opposing stance; one judge, 24 topics per arm) gives
Table~\ref{tab:authority-split} in Appendix~\ref{app:attacks}.

Every one-shot instruction complies with the neutrality request, and so does the
weight-resident margin adapter: each keeps the stance on at most one topic in twenty-four
(held 0--4.2\%). The margin adapter's collapse means the compiled arms' resistance (39.6\% base, 75.0\% with the engagement prompt)
belongs to the trained write function or its curriculum, not to weight-residence. Re-asserting
the persistence instruction before every turn beats them all (95.8\% held): revocability is a
property of \emph{one-shot} instructions, and a deployer who controls the serving loop can
exceed compiled-level resistance by re-injection, at the persistence instruction's usual costs
(17.5\% integrity-gate failure, 0\% concession to warranted evidence). Reversal is rare
everywhere (\(\le\)12.5\%), the bare host itself moves under this attack, and these are
single-judge cells over 24 topics.

Compilation has a second boundary, capacity: one merged state installs cleanly to about four
beliefs (Appendix~\ref{app:capacity}).\expref{app:exp-e2}%
\phantomsection\label{sec:boundary-bottleneck}

One reading of these results runs against the framing above, and belongs here rather than in
a rebuttal:
complying with a plausible operator correction is \emph{correct} behavior when the correction is
genuine. The attack is an attack only because the authority is claimed rather than real, and
nothing in the conversation distinguishes the two. Held rates in this table therefore measure
resistance to \emph{unverified} authority; a deployment that wants legitimate operator control
must provide it through a channel the interlocutor cannot imitate, which is an argument for
out-of-band control surfaces generally rather than for maximal in-conversation stubbornness.
Resistance to a false authority and obedience to a real one cannot be tuned separately: an
intervention that resists one resists the other.
